% Problems, Chapter 8: Sturm-Liouville Theory
% Arfken, Weber & Harris, Mathematical Methods for Physicists, 7th ed.
% Requires: amsmath, amssymb, tikz

\section*{Problems}

\Needspace{9\baselineskip}
\subsection*{8.2 Hermitian Operators}

\Needspace{5\baselineskip}
\par\bigskip
\noindent\textbf{8.2.1}\quad Show that Laguerre's ODE, Table 7.1, may be put
into self-adjoint form by multiplying by $e^{-x}$ and that $w(x)=e^{-x}$ is
the weighting function.
\solution
Write the equation as $xy''+(1-x)y'+ny=0$.  Multiplication by
$w=e^{-x}$ gives
\[
xe^{-x}y''+(1-x)e^{-x}y'+ne^{-x}y=0.
\]
The product rule gives $(xe^{-x})'=(1-x)e^{-x}$, so the first two
terms combine without changing the eigenvalue:
\[
\boxed{(xe^{-x}y')'+ne^{-x}y=0},\qquad
p(x)=xe^{-x},\quad w(x)=e^{-x}.
\]
Thus $A y=-w^{-1}(py')'=-xy''-(1-x)y'$ has eigenvalue $n$.
For functions in a common domain, integrate once in each scalar product:
\begin{align*}
\langle u,Av\rangle_w
&=-\int_a^b u^*(pv')'\,dx
=-[u^*pv']_a^b+\int_a^b p\,u'^*v'\,dx,\\
\langle Au,v\rangle_w
&=-\int_a^b(pu'^*)'v\,dx
=-[pu'^*v]_a^b+\int_a^b p\,u'^*v'\,dx.
\end{align*}
Subtracting cancels the identical integrals and gives
\[
\langle u,Av\rangle_w-\langle Au,v\rangle_w
=-\bigl[p(u^*v'-u'^*v)\bigr]_a^b,\qquad
\langle u,v\rangle_w=\int_a^b u^*v\,e^{-x}\,dx.
\]
This is the required weighted self-adjoint differential form; vanishing
of the displayed boundary form makes the operator Hermitian on that domain.
For polynomial solutions on $[0,\infty)$, both endpoint terms vanish,
as checked in Exercise 8.2.4.


\Needspace{5\baselineskip}
\par\bigskip
\noindent\textbf{8.2.2}\quad Show that the Hermite ODE, Table 7.1, may be put
into self-adjoint form by multiplying by $e^{-x^{2}}$ and that this gives
$w(x)=e^{-x^{2}}$ as the appropriate weighting function.
\solution
Hermite's equation is $y''-2xy'+2ny=0$.  Since
\[
\frac{d}{dx}e^{-x^2}=-2xe^{-x^2},
\]
multiplication by $e^{-x^2}$ gives
\[
e^{-x^2}y''-2xe^{-x^2}y'+2ne^{-x^2}y
=(e^{-x^2}y')'+2ne^{-x^2}y=0.
\]
Therefore
\[
\boxed{p=w=e^{-x^2},\qquad A=-w^{-1}(p\,d/dx)'=-d^2/dx^2+2x\,d/dx},
\]
where $Ay=2ny$.  In the scalar product
$\langle u,v\rangle_w=\int u^*v e^{-x^2}\,dx$, two integrations by
parts leave only $-[e^{-x^2}(u^*v'-u'^*v)]_a^b$.
It vanishes on $(-\infty,\infty)$ for polynomial solutions and their
linear combinations, establishing the claimed weight and Hermiticity.


\Needspace{5\baselineskip}
\par\bigskip
\noindent\textbf{8.2.3}\quad Show that the Chebyshev ODE, Table 7.1, may be
put into self-adjoint form by multiplying by $(1-x^{2})^{-1/2}$ and that this
gives $w(x)=(1-x^{2})^{-1/2}$ as the appropriate weighting function.
\solution
On $-1<x<1$, multiply $(1-x^2)y''-xy'+n^2y=0$ by
$(1-x^2)^{-1/2}$.  Because
\[
\frac{d}{dx}\sqrt{1-x^2}=-\frac{x}{\sqrt{1-x^2}},
\]
the resulting equation is
\[
\boxed{\left(\sqrt{1-x^2}\,y'\right)'
       +\frac{n^2}{\sqrt{1-x^2}}y=0}.
\]
Thus $p=\sqrt{1-x^2}$, $w=(1-x^2)^{-1/2}$ and
$A=-w^{-1}(p\,d/dx)'$ has eigenvalue $n^2$.
The scalar product is
\[
\boxed{\langle u,v\rangle_w=
\int_{-1}^{1}\frac{u^*(x)v(x)}{\sqrt{1-x^2}}\,dx}.
\]
The singularities of this weight are integrable: near either endpoint
they have the form a constant times $t^{-1/2}$, whose integral is
$2\sqrt t$.  Polynomial solutions therefore have finite norm, and
$p$ times any product of their values and first derivatives vanishes
at either endpoint.


\Needspace{5\baselineskip}
\par\bigskip
\noindent\textbf{8.2.4}\quad The Legendre, Chebyshev, Hermite, and Laguerre
equations, given in Table 7.1, have solutions that are polynomials. Show that
ranges of integration that guarantee that the Hermitian operator boundary
conditions will be satisfied are
\begin{enumerate}
\item[(a)] Legendre $[-1,1]$,\quad(b) Chebyshev $[-1,1]$,
\item[(c)] Hermite $(-\infty,\infty)$,\quad(d) Laguerre $[0,\infty)$.
\end{enumerate}
\solution
For $(pu')'+\lambda wu=0$, the boundary condition needed in the
orthogonality proof is
\[
\bigl[p(u^*v'-u'^*v)\bigr]_a^b=0.
\]
If $u,v$ are polynomials, then $u^*v'-u'^*v$ is another polynomial
(on the real axis).  We can therefore check each endpoint explicitly.
\begin{enumerate}
\item[(a)] Legendre: $p=1-x^2$ and $w=1$.
The polynomial factor is bounded at $x=\pm1$, while
$1-x^2\to0$.  Each boundary term is zero, and the norm integral over
$[-1,1]$ is finite.
\item[(b)] Chebyshev: $p=\sqrt{1-x^2}\to0$ at $\pm1$.
Bounded polynomial values and derivatives again make each boundary
term zero.  The norm exists because $w=(1-x^2)^{-1/2}$ is integrable.
\item[(c)] Hermite: $p=w=e^{-x^2}$.  Every polynomial is bounded in
magnitude by $C(1+|x|^N)$ for suitable $C,N$, and
$|x|^Ne^{-x^2}\to0$ as $x\to\pm\infty$.
Thus both boundary terms vanish and the weighted norms converge.
\item[(d)] Laguerre: $p=xe^{-x}$ and $w=e^{-x}$.
At $0$, the factor $x$ multiplies a bounded polynomial; at $\infty$,
$x^{N+1}e^{-x}\to0$.  Both boundary terms vanish, and the weighted
norms on $[0,\infty)$ converge.
\end{enumerate}
These conclusions use the polynomial hypothesis.  The same intervals
alone do not make arbitrary singular solutions satisfy the boundary
conditions, as Exercises 8.2.6 and 8.2.7 demonstrate.


\Needspace{5\baselineskip}
\par\bigskip
\noindent\textbf{8.2.5}\quad The functions $u_{1}(x)$ and $u_{2}(x)$ are
eigenfunctions of the same Hermitian operator but for distinct eigenvalues
$\lambda_{1}$ and $\lambda_{2}$. Prove that $u_{1}(x)$ and $u_{2}(x)$ are
linearly independent.
\solution
Suppose $c_1u_1+c_2u_2=0$, where eigenfunctions are nonzero by
definition.  Apply $A$ and then subtract $\lambda_2$ times the original
relation:
\[
c_1\lambda_1u_1+c_2\lambda_2u_2=0,
\qquad c_1(\lambda_1-\lambda_2)u_1=0.
\]
Since $\lambda_1\ne\lambda_2$ and $u_1\ne0$, we obtain $c_1=0$.
The original relation then implies $c_2=0$.  Hence the functions are
linearly independent.  Hermiticity is sufficient but is not needed
for this particular conclusion; only linearity and distinct eigenvalues
were used.


\Needspace{5\baselineskip}
\par\bigskip
\noindent\textbf{8.2.6}\quad Given that
\[
P_{1}(x)=x\quad\text{and}\quad
Q_{0}(x)=\frac{1}{2}\ln\left(\frac{1+x}{1-x}\right)
\]
are solutions of Legendre's differential equation (Table 7.1) corresponding to
different eigenvalues:
\begin{enumerate}
\item[(a)] Evaluate their orthogonality integral
  \[
  \int_{-1}^{1}\frac{x}{2}\ln\left(\frac{1+x}{1-x}\right)\,dx.
  \]
\item[(b)] Explain why these two functions are not orthogonal, that is, why
  the proof of orthogonality does not apply.
\end{enumerate}
\solution
\textbf{(a)} Put $Q=Q_0=\tfrac12\log((1+x)/(1-x))$.
Its derivative is
\[
Q'=\frac12\left(\frac1{1+x}+\frac1{1-x}\right)
=\frac1{1-x^2}.
\]
Integrate by parts on a closed subinterval of $(-1,1)$, choosing the
antiderivative $(x^2-1)/2$ of $x$:
\[
\int xQ\,dx=\frac{x^2-1}{2}Q
-\frac12\int(x^2-1)\frac{dx}{1-x^2}
=\frac{x^2-1}{2}Q+\frac{x}{2}.
\]
The endpoint products vanish because $t\log t\to0$ as $t\downarrow0$.
Taking both improper limits therefore gives
\[
\boxed{\int_{-1}^{1}P_1Q_0\,dx=0+\frac12(1-(-1))=1}.
\]
\textbf{(b)} For $A=-[(1-x^2)d/dx]'$, the eigenvalues of $P_1=x$
and $Q$ are $2$ and $0$.  Their boundary form is not zero:
\[
(1-x^2)(P_1Q'-P_1'Q)
=x-(1-x^2)Q\longrightarrow
\begin{cases}1,&x\to1,\\-1,&x\to-1.\end{cases}
\]
Its endpoint difference is $2$.  Indeed, subtracting their differential
equations gives
\[
\bigl[(1-x^2)(P_1Q'-P_1'Q)\bigr]_{-1}^{1}
=(2-0)\int_{-1}^{1}P_1Q\,dx=2.
\]
Although $Q$ is square integrable, it is singular at the endpoints and
does not belong to the same boundary-condition domain as the Legendre
polynomials.  The Hermitian orthogonality theorem thus does not apply
to this pair.


\Needspace{5\baselineskip}
\par\bigskip
\noindent\textbf{8.2.7}\quad $T_{0}(x)=1$ and $V_{1}(x)=(1-x^{2})^{1/2}$ are
solutions of the Chebyshev differential equation corresponding to different
eigenvalues. Explain, in terms of the boundary conditions, why these two
functions are not orthogonal on the range $(-1,1)$ with the weighting function
found in Exercise 8.2.3.
\solution
With $u=T_0=1$, $v=V_1=\sqrt{1-x^2}$, $p=\sqrt{1-x^2}$ and
$w=1/p$, the weighted inner product is
\[
\boxed{\langle u,v\rangle_w=\int_{-1}^{1}1\,dx=2\ne0}.
\]
Here $u'=0$, $v'=-x/p$, so
\[
p(uv'-u'v)=-x,
\qquad \bigl[p(uv'-u'v)\bigr]_{-1}^{1}=-2.
\]
The two eigenvalues are $0$ and $1$: explicitly,
$(pu')'=0$ and $(pv')'=-1=-wv$.
Consequently the boundary identity gives
$-2=(0-1)\langle u,v\rangle_w$, consistently with the integral.
The divergence of $v'$ cancels the zero of $p$; the boundary form
does not vanish even though $v$ itself vanishes at both endpoints.
They are therefore not eigenfunctions in a common Hermitian domain
containing the Chebyshev polynomials.


\Needspace{5\baselineskip}
\par\bigskip
\noindent\textbf{8.2.8}\quad A set of functions $u_{n}(x)$ satisfies the
Sturm-Liouville equation
\[
\frac{d}{dx}\left[p(x)\frac{d}{dx}u_{n}(x)\right]+\lambda_{n}w(x)u_{n}(x)=0.
\]
The functions $u_{m}(x)$ and $u_{n}(x)$ satisfy boundary conditions that lead
to orthogonality. The corresponding eigenvalues $\lambda_{m}$ and
$\lambda_{n}$ are distinct. Prove that for appropriate boundary conditions,
$u'_{m}(x)$ and $u'_{n}(x)$ are orthogonal with $p(x)$ as a weighting
function.
\solution
Multiply the equation for $u_n$ by $u_m^*$ and integrate:
\[
0=\int_a^b u_m^*(pu_n')'\,dx
 +\lambda_n\int_a^b wu_m^*u_n\,dx.
\]
One integration by parts gives the exact identity
\[
\boxed{\int_a^b p\,u_m'^*u_n'\,dx
=\bigl[pu_m^*u_n'\bigr]_a^b
 +\lambda_n\int_a^b wu_m^*u_n\,dx}.
\]
For distinct eigenvalues the last integral is zero by the assumed
orthogonality.  Thus the additional sufficient boundary condition is
\[
\bigl[pu_m^*u_n'\bigr]_a^b=0,
\]
which yields $\int_a^b p\,u_m'^*u_n'\,dx=0$.
This holds, for example, under common Dirichlet conditions, common
Neumann conditions, or a common choice of Dirichlet or Neumann
condition at each endpoint.  It also holds for periodic values of $u$
and $pu'$, when the endpoint products cancel, or when $p\to0$ and
the products $u_m^*u_n'$ remain bounded.

The vanishing of the antisymmetrized boundary form alone need not
suffice.  Under real Robin conditions $u'(a)=h_a u(a)$ and
$u'(b)=h_bu(b)$ that form vanishes, but the boundary term here is
\[
p(b)h_bu_m^*(b)u_n(b)-p(a)h_au_m^*(a)u_n(a),
\]
which requires a separate check.  For normalized eigenfunctions and
the stronger vanishing condition, the same calculation with $m=n$
also gives $\int_a^b p|u_n'|^2\,dx=\lambda_n$.


\Needspace{5\baselineskip}
\par\bigskip
\noindent\textbf{8.2.9}\quad Linear operator $\mathcal{A}$ has $n$ distinct
eigenvalues and $n$ corresponding eigenfunctions:
$\mathcal{A}\psi_{i}=\lambda_{i}\psi_{i}$. Show that the $n$ eigenfunctions
are linearly independent. Do not assume $\mathcal{A}$ to be Hermitian.

\emph{Hint.}~Assume linear dependence, i.e., that
$\psi_{n}=\sum_{i=1}^{n-1}a_{i}\psi_{i}$. Use this relation and the
operator-eigenfunction equation first in one order and then in the reverse
order. Show that a contradiction results.
\solution
We prove the statement by induction, so that comparison of coefficients
will not presuppose the independence we are trying to establish.
One nonzero eigenfunction is independent.  Suppose any $n-1$
eigenfunctions with distinct eigenvalues are independent, and consider
\[
\sum_{i=1}^n c_i\psi_i=0.
\]
Apply $\mathcal A-\lambda_n I$ to this equality.  Linearity and the
eigenfunction equations give
\[
\sum_{i=1}^{n-1}c_i(\lambda_i-\lambda_n)\psi_i=0.
\]
By the induction hypothesis every coefficient in this shorter sum is
zero.  Since $\lambda_i-\lambda_n\ne0$, $c_i=0$ for $i<n$.
The original equality becomes $c_n\psi_n=0$, hence $c_n=0$.
This completes the induction without Hermiticity and without division
by any individual eigenvalue (which may be zero).


\Needspace{5\baselineskip}
\par\bigskip
\noindent\textbf{8.2.10}\quad The ultraspherical polynomials
$C_{n}^{(\alpha)}(x)$ are solutions of the differential equation
\[
\left\{(1-x^{2})\frac{d^{2}}{dx^{2}}-(2\alpha+1)x\frac{d}{dx}+n(n+2\alpha)
\right\}C_{n}^{(\alpha)}(x)=0.
\]
\begin{enumerate}
\item[(a)] Transform this differential equation into self-adjoint form.
\item[(b)] Find an interval of integration and weighting factor that make
  $C_{n}^{(\alpha)}(x)$ of the same $\alpha$ but different $n$ orthogonal.
\end{enumerate}
\emph{Note.}~Assume that your solutions are polynomials.
\erratum{Take real $\alpha>-1/2$. At $\alpha=0$, use nonzero
limiting normalizations for degrees $n\ge1$.}
\solution
\erratum{For a positive integrable weight supporting all degrees,
assume real $\alpha>-\tfrac12$.  Eigenfunctions must be nonzero;
at $\alpha=0$ use the nonzero Chebyshev limits rather than the
identically zero standard $C_n^{(0)}$ for $n\ge1$.}

\textbf{(a)} If $w$ is the multiplier and $p=(1-x^2)w$, the
coefficient of $y'$ must be $p'=-(2\alpha+1)xw$.  Expanding this
condition gives
\[
(1-x^2)w'-2xw=-(2\alpha+1)xw,
\qquad \frac{w'}w=\frac{(1-2\alpha)x}{1-x^2}.
\]
Since $d\log(1-x^2)/dx=-2x/(1-x^2)$, integration yields
\[
w=(1-x^2)^{\alpha-1/2},\qquad p=(1-x^2)^{\alpha+1/2},
\]
where the irrelevant positive constant has been set to one.  Hence
\[
\boxed{\bigl[(1-x^2)^{\alpha+1/2}y'\bigr]'
+n(n+2\alpha)(1-x^2)^{\alpha-1/2}y=0}.
\]
\textbf{(b)} On $(-1,1)$ this weight is positive.  At either endpoint
it behaves as a positive constant times $t^{\alpha-1/2}$, whose
integral is finite exactly when $\alpha>-1/2$.  The same inequality
makes $p\to0$, so the boundary form vanishes for polynomial solutions.
For $m\ne n$ in $\{0,1,2,\ldots\}$,
\[
\lambda_n-\lambda_m=(n-m)(n+m+2\alpha)\ne0,
\]
because $n+m\ge1$ and $2\alpha>-1$.  Subtraction of the two
self-adjoint equations therefore gives
\[
\boxed{\int_{-1}^{1}(1-x^2)^{\alpha-1/2}
 C_m^{(\alpha)}(x)C_n^{(\alpha)}(x)\,dx=0\quad(m\ne n)}.
\]
To explain the qualification at zero, the generating function
$(1-2xt+t^2)^{-\alpha}=\sum_{n\ge0}C_n^{(\alpha)}(x)t^n$
is $1$ when $\alpha=0$.  For $n\ge1$ its nonzero limiting normalization
is $\lim_{\alpha\to0}C_n^{(\alpha)}(x)/\alpha=2T_n(x)/n$.
To see the coefficient explicitly, differentiate the generating function
with respect to $\alpha$ at zero. Put $x=\cos\theta$ and factor
$1-2xt+t^2=(1-e^{i\theta}t)(1-e^{-i\theta}t)$. For $|t|<1$,
\begin{align*}
-\log(1-2xt+t^2)
&=\sum_{n\ge1}\frac{e^{in\theta}+e^{-in\theta}}n t^n\\
&=2\sum_{n\ge1}\frac{\cos(n\theta)}n t^n
=2\sum_{n\ge1}\frac{T_n(x)}n t^n.
\end{align*}
The resulting Chebyshev family has precisely the weight in
Exercise 8.2.3.


\Needspace{9\baselineskip}
\subsection*{8.3 ODE Eigenvalue Problems}

\Needspace{5\baselineskip}
\par\bigskip
\noindent\textbf{8.3.1}\quad Solve the Legendre equation
\[
(1-x^{2})y''-2xy'+n(n+1)y=0
\]
by direct series substitution.
\begin{enumerate}
\item[(a)] Verify that the indicial equation is
  \[
  s(s-1)=0.
  \]
\item[(b)] Using $s=0$ and setting the coefficient $a_{1}=0$, obtain a series
  of even powers of $x$:
  \[
  y_{\mathrm{even}}=a_{0}\left[1-\frac{n(n+1)}{2!}x^{2}
  +\frac{(n-2)n(n+1)(n+3)}{4!}x^{4}+\cdots\right],
  \]
  where
  \[
  a_{j+2}=\frac{j(j+1)-n(n+1)}{(j+1)(j+2)}a_{j}.
  \]
\item[(c)] Using $s=1$ and noting that the coefficient $a_{1}$ must be zero,
  develop a series of odd powers of $x$:
  \[
  y_{\mathrm{odd}}=a_{0}\left[x-\frac{(n-1)(n+2)}{3!}x^{3}
  +\frac{(n-3)(n-1)(n+2)(n+4)}{5!}x^{5}+\cdots\right],
  \]
  where
  \[
  a_{j+2}=\frac{(j+1)(j+2)-n(n+1)}{(j+2)(j+3)}a_{j}.
  \]
\item[(d)] Show that both solutions, $y_{\mathrm{even}}$ and
  $y_{\mathrm{odd}}$, diverge for $x=\pm1$ if the series continue to infinity.
  (Compare with Exercise 1.2.5.)
\item[(e)] Finally, show that by an appropriate choice of $n$, one series at a
  time may be converted into a polynomial, thereby avoiding the divergence
  catastrophe. In quantum mechanics this restriction of $n$ to integral values
  corresponds to quantization of angular momentum.
\end{enumerate}
\solution
\textbf{(a)} Substitute $y=\sum_{j\ge0}a_jx^{j+s}$, with $a_0\ne0$.
Before combining powers, the equation reads
\[
\sum_{j\ge0}(j+s)(j+s-1)a_jx^{j+s-2}
+\sum_{j\ge0}\{n(n+1)-(j+s)(j+s+1)\}a_jx^{j+s}=0.
\]
The coefficient of $x^{s-2}$ is $s(s-1)a_0$, and that of $x^{s-1}$
is $s(s+1)a_1$.  Thus $s=0$ or $1$; $a_1$ is free for $s=0$
and is zero for $s=1$.  Matching the remaining powers gives
\[
\boxed{a_{j+2}=\frac{(j+s-n)(j+s+n+1)}{(j+s+1)(j+s+2)}a_j}.
\]
The point $x=0$ is ordinary; the two values of $s$ simply select the
even and odd Taylor solutions.

\textbf{(b)} For $s=0$, set $a_1=0$.  Then
\[
a_2=-\frac{n(n+1)}{2}a_0,\qquad
 a_4=\frac{(2-n)(n+3)}{3\cdot4}a_2
=\frac{(n-2)n(n+1)(n+3)}{4!}a_0.
\]
In general, with an empty product equal to one,
\[
y_e=a_0\sum_{k=0}^{\infty}
\frac{\prod_{r=0}^{k-1}(2r-n)(2r+n+1)}{(2k)!}x^{2k}.
\]
This gives the stated even series and its recurrence.

\textbf{(c)} For $s=1$, $a_1=0$ is forced by $2a_1=0$.  Now
\[
a_2=-\frac{(n-1)(n+2)}{3!}a_0,
\qquad
 a_4=\frac{(3-n)(n+4)}{4\cdot5}a_2
=\frac{(n-3)(n-1)(n+2)(n+4)}{5!}a_0.
\]
Thus the independent odd solution is
\[
y_o=a_0\sum_{k=0}^{\infty}
\frac{\prod_{r=0}^{k-1}(2r+1-n)(2r+n+2)}{(2k+1)!}x^{2k+1}.
\]
The two occurrences of $a_0$ are independent arbitrary constants.

\textbf{(d)} Let $c_k$ be the coefficient of $x^{2k+s}$ in either
nonterminating branch, with $s=0$ or $1$.  The exact recurrence implies
\[
\frac{c_{k+1}}{c_k}
=\frac{(2k+s)(2k+s+1)-n(n+1)}{(2k+s+1)(2k+s+2)}
=1-\frac{2(2k+s+1)+n(n+1)}{(2k+s+1)(2k+s+2)}.
\]
The subtracted numerator is $4k+O(1)$ and the denominator is
$4k^2+O(k)$, so the ratio is $1-1/k+O(k^{-2})$.
Consequently
\[
\frac{(k+1)c_{k+1}}{kc_k}=1+O(k^{-2}).
\]
Starting beyond all possible zeros, let $b_k=kc_k$. The recurrence
for $b_k$ is $b_{k+1}/b_k=1+O(k^{-2})$. The logarithms of these
ratios have an absolutely convergent sum, so their product has a
finite nonzero limit $C$. Its tail has logarithm
$O(\sum_{j\ge k}j^{-2})=O(k^{-1})$; hence
$b_k=C[1+O(k^{-1})]$ and
\[
c_k=\frac Ck+O(k^{-2}),\qquad C\ne0.
\]
At $x=1$, the $C/k$ terms have divergent partial sums, while the
$O(k^{-2})$ remainder is absolutely summable. At $x=-1$, every
term acquires the same factor $(-1)^s$,
so it diverges there too.  The ratio test gives convergence for $|x|<1$.
\erratum{The manual's observation that $c_{k+1}/c_k\to1$ is not by
itself a divergence proof; the $1/k$ asymptotic above supplies the
necessary endpoint argument.}

\textbf{(e)} A branch terminates when an exponent $r=2k+s\ge0$
satisfies $(r-n)(r+n+1)=0$.  With the conventional representative
$n\ge0$, this means $n=r$: even nonnegative $n$ truncates the even
branch, and odd positive $n$ truncates the odd branch.  The other
branch remains nonterminating and singular at the endpoints.
\erratum{If negative parameters are allowed, $n=-r-1$ produces the
same eigenvalue $n(n+1)=r(r+1)$ and the same polynomial; restricting
to $n\ge0$ removes this duplicate labeling.}
Normalizing the surviving polynomial by $P_n(1)=1$ gives the Legendre
polynomial, for example $1,x,(3x^2-1)/2,(5x^3-3x)/2$ for $n=0,1,2,3$.


\Needspace{5\baselineskip}
\par\bigskip
\noindent\textbf{8.3.2}\quad Show that with the weight factor $\exp(-x^{2})$
and the interval $-\infty<x<\infty$ for the scalar product, the Hermite ODE
eigenvalue problem is Hermitian.
\erratum{Use a common domain with vanishing boundary form at both
infinities, such as the span of the polynomial eigenfunctions.}
\solution
\erratum{A common domain satisfying the endpoint condition is required
in addition to the weight and interval.  The polynomial eigenfunctions
satisfy it.}
Write the eigenvalue problem as $Ay=2\alpha y$, where
\[
A=-\frac{d^2}{dx^2}+2x\frac{d}{dx}
=-e^{x^2}\frac{d}{dx}\left(e^{-x^2}\frac{d}{dx}\right).
\]
For finite endpoints $a,b$, integration by parts yields
\begin{align*}
\langle u,Av\rangle_w
&=-\int_a^b u^*(e^{-x^2}v')'\,dx\\
&=-[u^*e^{-x^2}v']_a^b
  +\int_a^b u'^*e^{-x^2}v'\,dx,\\
\langle Au,v\rangle_w
&=-[u'^*e^{-x^2}v]_a^b
  +\int_a^b u'^*e^{-x^2}v'\,dx.
\end{align*}
Their difference is
\[
-\bigl[e^{-x^2}(u^*v'-u'^*v)\bigr]_a^b.
\]
For polynomials $u,v$, the Gaussian defeats every power of $|x|$,
so both endpoint limits are zero as $a\to-\infty$, $b\to\infty$.
All displayed integrals converge.  The identity
$\langle u,Av\rangle_w=\langle Au,v\rangle_w$ follows on the
invariant domain of finite linear combinations of these polynomials.
More generally it holds on any common operator domain where the
integrals exist and the same boundary form vanishes.

\erratum{Weighted square integrability alone, as asserted in the
manual, is insufficient to establish this boundary condition.}
For example $u=1$ and $v=\sin(e^{x^2})$ are both in
$L^2(\mathbb R,e^{-x^2}dx)$, since they are bounded, but
\[
e^{-x^2}(uv'-u'v)=2x\cos(e^{x^2})
\]
has no zero limit: at $x_k=\sqrt{\log(2\pi k)}$ it equals $2x_k$.
Such a function is excluded from the domain used in the Hermitian
problem; square integrability does not replace an operator-domain
condition.


\Needspace{5\baselineskip}
\par\bigskip
\noindent\textbf{8.3.3}\quad
\begin{enumerate}
\item[(a)] Develop series solutions for Hermite's differential equation
  \[
  y''-2xy'+2\alpha y=0.
  \]
  \par\nobreak\hfill\emph{ANS.}\enspace $s(s-1)=0$, indicial equation.\par
  For $s=0$,
  \[
  a_{j+2}=2a_{j}\frac{j-\alpha}{(j+1)(j+2)}\quad(j\text{ even}),
  \]
  \[
  y_{\mathrm{even}}=a_{0}\left[1+\frac{2(-\alpha)x^{2}}{2!}
  +\frac{2^{2}(-\alpha)(2-\alpha)x^{4}}{4!}+\cdots\right].
  \]
  For $s=1$,
  \[
  a_{j+2}=2a_{j}\frac{j+1-\alpha}{(j+2)(j+3)}\quad(j\text{ even}),
  \]
  \[
  y_{\mathrm{odd}}=a_{1}\left[x+\frac{2(1-\alpha)x^{3}}{3!}
  +\frac{2^{2}(1-\alpha)(3-\alpha)x^{5}}{5!}+\cdots\right].
  \]
\item[(b)] Show that both series solutions are convergent for all $x$, the
  ratio of successive coefficients behaving, for a large index, like the
  corresponding ratio in the expansion of $\exp(x^{2})$.
\item[(c)] Show that by appropriate choice of $\alpha$, the series solutions
  may be cut off and converted to finite polynomials. (These polynomials,
  properly normalized, become the Hermite polynomials in Section 18.1.)
\end{enumerate}
\solution
\textbf{(a)} Substitution of $y=\sum_{j\ge0}a_jx^{j+s}$ gives
\[
\sum_{j\ge0}(j+s)(j+s-1)a_jx^{j+s-2}
+\sum_{j\ge0}2(\alpha-j-s)a_jx^{j+s}=0.
\]
Thus $s(s-1)=0$, $s(s+1)a_1=0$, and
\[
\boxed{a_{j+2}=\frac{2(j+s-\alpha)}{(j+s+1)(j+s+2)}a_j}.
\]
For the even solution take $s=0$, $a_1=0$, $a_0=A$.  Successive
substitutions give
\[
a_2=-\alpha A,\qquad
 a_4=\frac{2(2-\alpha)}{3\cdot4}(-\alpha A),
\]
and hence
\[
y_e=A\sum_{k\ge0}
\frac{2^k\prod_{r=0}^{k-1}(2r-\alpha)}{(2k)!}x^{2k}
=A\left[1-\alpha x^2+\frac{\alpha(\alpha-2)}6x^4+\cdots\right].
\]
For the odd solution take $s=1$.  Its Frobenius coefficient $a_1$
is forced to zero; call its free $a_0$ coefficient $B$, which is the
coefficient of $x$ (denoted $a_1$ in the printed answer).  Then
\[
a_2=\frac{2(1-\alpha)}{2\cdot3}B,\qquad
 a_4=\frac{2(3-\alpha)}{4\cdot5}a_2,
\]
so
\[
y_o=B\sum_{k\ge0}
\frac{2^k\prod_{r=0}^{k-1}(2r+1-\alpha)}{(2k+1)!}x^{2k+1}.
\]
Its first terms are exactly $B[x+2(1-\alpha)x^3/3!
+2^2(1-\alpha)(3-\alpha)x^5/5!+\cdots]$.

\textbf{(b)} For a fixed $x$, the ratio of successive nonzero
terms of either parity has magnitude
\[
|x|^2\left|\frac{2(j+s-\alpha)}{(j+s+1)(j+s+2)}\right|
\sim\frac{2|x|^2}{j}\longrightarrow0.
\]
Both series therefore converge absolutely for every finite real or
complex $x$.  With $j=2k$, the ratio is asymptotic to $x^2/k$,
which is the large-$k$ behavior of the ratio $x^2/(k+1)$ for
$e^{x^2}=\sum_{k\ge0}x^{2k}/k!$.

\textbf{(c)} Termination requires $j+s=\alpha$ at a nonnegative
integer exponent.  Thus $\alpha=N\in\{0,1,2,\ldots\}$ gives one
polynomial: the even branch for even $N$ and the odd branch for odd
$N$.  The other branch does not terminate, and its coefficient must
be set to zero if a polynomial is required.  Choosing leading
coefficient $2^N$ yields $H_N$, beginning with
\[
H_0=1,\quad H_1=2x,\quad H_2=4x^2-2,\quad H_3=8x^3-12x.
\]


\Needspace{5\baselineskip}
\par\bigskip
\noindent\textbf{8.3.4}\quad Laguerre's ODE is
\[
xL''_{n}(x)+(1-x)L'_{n}(x)+nL_{n}(x)=0.
\]
Develop a series solution and select the parameter $n$ to make your series a
polynomial.
\solution
Use $y=\sum_{j\ge0}a_jx^{j+s}$.  The $xy''+y'$ terms combine as
\[
xy''+y'=\sum_{j\ge0}(j+s)^2a_jx^{j+s-1},
\]
whereas $-xy'+ny=\sum_{j\ge0}(n-j-s)a_jx^{j+s}$.
The lowest coefficient is $s^2a_0$, so the repeated indicial root
is $s=0$.  Matching $x^{j+s}$ gives
\[
(j+s+1)^2a_{j+1}+(n-j-s)a_j=0,
\qquad
\boxed{a_{j+1}=\frac{j-n}{(j+1)^2}a_j\quad(s=0)}.
\]
\erratum{The recurrence advances by one index; the manual's
$a_{j+2}$ on the left is a typographical error.}
Iteration yields
\[
a_1=-na_0,\quad a_2=\frac{n(n-1)}{(2!)^2}a_0,
\quad a_3=-\frac{n(n-1)(n-2)}{(3!)^2}a_0,
\]
and, in general,
\[
y=a_0\sum_{j=0}^{\infty}
\frac{(-1)^j n(n-1)\cdots(n-j+1)}{(j!)^2}x^j.
\]
For a nonterminating solution the term ratio is asymptotic to $x/j$,
so this regular solution is entire.  It becomes a polynomial precisely
when $n=N$ is a nonnegative integer; then $a_{N+1}=0$ and all later
coefficients vanish.  With $a_0=1$,
\[
\boxed{L_N(x)=\sum_{j=0}^{N}(-1)^j\binom Nj\frac{x^j}{j!}}.
\]
For example $L_0=1$, $L_1=1-x$, and $L_2=1-2x+x^2/2$.
There is only one analytic solution near zero up to scale, since
$a_0$ determines every coefficient.  Reduction of order supplies a
second local solution
\[
y_2=y_1\int\frac{e^x}{x\,y_1(x)^2}\,dx,
\]
where $y_1(0)=1$.  Its integrand is $1/x+O(1)$, so it contains
$\log x$ and cannot be a regular polynomial at zero.


\Needspace{5\baselineskip}
\par\bigskip
\noindent\textbf{8.3.5}\quad Solve the Chebyshev equation
\[
(1-x^{2})T''_{n}-xT'_{n}+n^{2}T_{n}=0,
\]
by series substitution. What restrictions are imposed on $n$ if you demand
that the series solution converge for $x=\pm1$?
\par\nobreak\hfill\emph{ANS.}\enspace
The infinite series does converge for $x=\pm1$ and no restriction on $n$
exists (compare with Exercise 1.2.6).
\solution
Substitute $y=\sum_{j\ge0}a_jx^{j+s}$.  The coefficients of
$x^{s-2}$ and $x^{s-1}$ give $s(s-1)=0$ and $s(s+1)a_1=0$.
The other coefficients give
\[
(j+s+1)(j+s+2)a_{j+2}+[n^2-(j+s)^2]a_j=0,
\]
so
\[
\boxed{a_{j+2}=\frac{(j+s)^2-n^2}{(j+s+1)(j+s+2)}a_j}.
\]
The independent even and odd solutions, normalized to $y_e(0)=1$
and $y_o'(0)=1$, are therefore
\begin{align*}
y_e&=\sum_{k\ge0}
\frac{\prod_{r=0}^{k-1}[(2r)^2-n^2]}{(2k)!}x^{2k}
=1-\frac{n^2x^2}{2!}+\frac{n^2(n^2-4)x^4}{4!}+\cdots,\\
y_o&=\sum_{k\ge0}
\frac{\prod_{r=0}^{k-1}[(2r+1)^2-n^2]}{(2k+1)!}x^{2k+1}
=x+\frac{(1-n^2)x^3}{3!}
 +\frac{(1-n^2)(9-n^2)x^5}{5!}+\cdots.
\end{align*}
For a nonterminating branch put $c_k=a_{2k}$ in its own Frobenius
notation.  Its ratio is
\[
\frac{c_{k+1}}{c_k}
=\frac{(2k+s)^2-n^2}{(2k+s+1)(2k+s+2)}
=1-\frac{3(2k+s)+2+n^2}{(2k+s+1)(2k+s+2)}.
\]
The fraction subtracted from one is
$(6k+O(1))/(4k^2+O(k))=3/(2k)+O(k^{-2})$.
Multiplying $c_k$ by $k^{3/2}$ makes the new successive ratio
$1+O(k^{-2})$, whose product converges to a nonzero constant.
Thus $c_k\sim Ck^{-3/2}$.  Comparison with $\sum k^{-3/2}$ proves
absolute convergence at both $x=1$ and $x=-1$.
Terminating branches also converge there.  Hence
\[
\boxed{\text{endpoint convergence imposes no restriction on }n.}
\]
For an independent check, set $x=\sin t$ with $|t|<\pi/2$.
Writing $Y(t)=y(\sin t)$, the chain rule gives
\[
Y'=y'\cos t,\qquad
Y''=y''\cos^2t-y'\sin t=(1-x^2)y''-xy'.
\]
Thus $Y''+n^2Y=0$. The initial conditions at zero select
\[
y_e=\cos(n\arcsin x),\qquad
y_o=\begin{cases}\sin(n\arcsin x)/n,&n\ne0,\\\arcsin x,&n=0.
\end{cases}
\]
Their endpoint values are finite.  Polynomial solutions instead
require $n^2=N^2$, $N$ a nonnegative integer, and selection of the
branch of parity $N$.  Endpoint convergence and polynomial
termination are therefore different requirements.


\Needspace{5\baselineskip}
\par\bigskip
\noindent\textbf{8.3.6}\quad Solve
\[
(1-x^{2})U''_{n}(x)-3xU'_{n}(x)+n(n+2)U_{n}(x)=0,
\]
choosing the root of the indicial equation to obtain a series of odd powers of
$x$. Since the \erratum{nonterminating} series will diverge for $x=1$, choose $n$ to convert it into a
polynomial.
\solution
Substitution of $y=\sum_{j\ge0}a_jx^{j+s}$ gives
$s(s-1)=0$, $s(s+1)a_1=0$, and
\[
a_{j+2}
=\frac{(j+s)(j+s+2)-n(n+2)}{(j+s+1)(j+s+2)}a_j.
\]
For the odd branch choose $s=1$, which forces $a_1=0$.  Factor the
numerator to obtain
\[
\boxed{a_{j+2}=\frac{(j+1-n)(j+n+3)}{(j+2)(j+3)}a_j
\quad(j=0,2,4,\ldots)}.
\]
In particular,
\[
a_2=-\frac{(n-1)(n+3)}{3!}a_0,\qquad
 a_4=\frac{(n-3)(n-1)(n+3)(n+5)}{5!}a_0,
\]
so
\[
y=a_0\left[x-\frac{(n-1)(n+3)}{3!}x^3
+\frac{(n-3)(n-1)(n+3)(n+5)}{5!}x^5+\cdots\right].
\]
For the nonterminating odd series, $c_k=a_{2k}$ satisfies
\[
\frac{c_{k+1}}{c_k}
=\frac{(2k+1)(2k+3)-n(n+2)}{(2k+2)(2k+3)}
=1-\frac{2k+3+n(n+2)}{(2k+2)(2k+3)}.
\]
The subtracted fraction is $(2k+O(1))/(4k^2+O(k))
=1/(2k)+O(k^{-2})$.
As in Exercise 8.3.1, a convergent product shows
$c_k\sim Ck^{-1/2}$, $C\ne0$; the terms have a fixed eventual sign
for real $n$, so the series diverges at $x=1$ (and at $-1$).
Termination occurs at an odd exponent $N=2k+1$ when
\[
(N-n)(N+n+2)=0.
\]
With the usual convention $n\ge0$, choose $n=N=1,3,5,\ldots$;
then the degree is $N$.  \erratum{Allowing negative labels also gives
$n=-N-2=-3,-5,\ldots$, which produce the same eigenvalue and
polynomial.}
The standard second-kind Chebyshev normalization is
$a_0=(-1)^{(N-1)/2}(N+1)$.  The first three polynomials are
\[
U_1=2x,\qquad U_3=8x^3-4x,\qquad U_5=32x^5-32x^3+6x.
\]
This normalization can be checked from
$U_N(\cos\theta)=\sin((N+1)\theta)/\sin\theta$. Since
$d/dx=-(\sin\theta)^{-1}d/d\theta$, differentiation at
$x=0$ ($\theta=\pi/2$) gives
\[
a_0=U_N'(0)=-(N+1)\cos\frac{(N+1)\pi}{2}
=(-1)^{(N-1)/2}(N+1),
\]
where the last equality uses odd $N$.


\Needspace{9\baselineskip}
\subsection*{8.4 Variation Method}

\Needspace{5\baselineskip}
\par\bigskip
\noindent\textbf{8.4.1}\quad A function that is normalized on the interval
$0\le x<\infty$ with an unweighted scalar product is
\[
\psi=2\alpha^{3/2}xe^{-\alpha x},\qquad
\erratum{\alpha>0\ \text{real}}.
\]
\begin{enumerate}
\item[(a)] Verify the normalization.
\item[(b)] Verify that for this $\psi$, $\langle x^{-1}\rangle=\alpha$.
\item[(c)] Verify that for this $\psi$,
  $\langle d^{2}/dx^{2}\rangle=-\alpha^{2}$.
\item[(d)] Use the variation method to find the value of $\alpha$ that
  minimizes
  \[
  \left\langle\psi\left|-\frac{1}{2}\frac{d^{2}}{dx^{2}}-\frac{1}{x}\right|
  \psi\right\rangle,
  \]
  and find the minimum value of this expectation value.
\end{enumerate}
\solution
\erratum{Take $\alpha>0$ real; this is necessary for the stated
normalized trial function.}
For $\beta>0$, define $I_k(\beta)=\int_0^\infty x^ke^{-\beta x}\,dx$.
Direct integration gives $I_0=1/\beta$, and integration by parts gives
\[
I_k=\left[-\frac{x^ke^{-\beta x}}\beta\right]_0^\infty
+\frac{k}\beta I_{k-1}=\frac{k}\beta I_{k-1}\quad(k\ge1).
\]
Therefore $I_k=k!/\beta^{k+1}$.

\textbf{(a)} The normalization integral is
\[
\int_0^\infty|\psi|^2\,dx
=4\alpha^3 I_2(2\alpha)
=4\alpha^3\frac{2!}{(2\alpha)^3}=1.
\]
\textbf{(b)} Multiplication by $x^{-1}$ removes one power of $x$:
\[
\boxed{\langle x^{-1}\rangle
=4\alpha^3 I_1(2\alpha)
=4\alpha^3\frac{1!}{(2\alpha)^2}=\alpha}.
\]
The apparent singularity at $0$ is harmless because the integrand is
proportional to $x e^{-2\alpha x}$.

\textbf{(c)} Differentiate with the exponential retained:
\[
\psi'=2\alpha^{3/2}(1-\alpha x)e^{-\alpha x},\qquad
\psi''=2\alpha^{3/2}(\alpha^2x-2\alpha)e^{-\alpha x}.
\]
\erratum{The factor $e^{-\alpha x}$ is missing from the manual's
displayed second derivative, although its following integral is correct.}
Consequently
\begin{align*}
\langle d^2/dx^2\rangle
&=4\alpha^3\int_0^\infty
 (\alpha^2x^2-2\alpha x)e^{-2\alpha x}\,dx\\
&=4\alpha^3\left[\alpha^2\frac{2!}{(2\alpha)^3}
 -2\alpha\frac{1!}{(2\alpha)^2}\right]
=\boxed{-\alpha^2}.
\end{align*}
The sign also follows by integration by parts:
$\int\psi\psi''=[\psi\psi']_0^\infty-\int(\psi')^2
=-\int(\psi')^2$, since $\psi(0)=0$ and the exponential removes
the boundary term at infinity.

\textbf{(d)} Combine the two expectation values:
\[
E(\alpha)=-\tfrac12(-\alpha^2)-\alpha
=\frac{\alpha^2}{2}-\alpha
=\frac{(\alpha-1)^2}{2}-\frac12.
\]
Thus the unique minimum for $\alpha>0$ is
\[
\boxed{\alpha_{\mathrm{opt}}=1,\qquad E_{\min}=-\tfrac12}.
\]
Equivalently $E'=\alpha-1=0$ and $E''=1>0$.
As a direct check, at $\alpha=1$ the expression for $\psi''$ gives
$(-\tfrac12d^2/dx^2-1/x)\psi=-\tfrac12\psi$; the optimized trial
function is an exact eigenfunction in this case.
